\documentclass[11pt]{article}

\usepackage[margin=1in]{geometry}
\usepackage{amsmath,amssymb,mathtools,bm}
\usepackage{booktabs}
\usepackage{microtype}
\usepackage{siunitx}
\usepackage{hyperref}
\usepackage[nameinlink,noabbrev]{cleveref}
\usepackage{xcolor}
\usepackage{enumitem}

\hypersetup{colorlinks=true,linkcolor=blue!45!black,citecolor=blue!45!black,urlcolor=blue!45!black}
\setlist[itemize]{topsep=3pt,itemsep=2pt}
\setlist[enumerate]{topsep=3pt,itemsep=2pt}

\newcommand{\dd}{\,\mathrm d}
\newcommand{\rhobar}{\bar{\rho}_{m,0}}
\newcommand{\Mpc}{\mathrm{Mpc}}
\newcommand{\Msun}{M_\odot}
\newcommand{\LCDM}{\Lambda\mathrm{CDM}}

\title{Research Note\\\large From the Linear Matter Power Spectrum to the Halo Mass Function}
\author{Blake Bird}
\date{}


\begin{document}
\maketitle

\begin{abstract}
This note establishes the independent physics and numerical layer required before studying Early Dark Energy (EDE) effects on halo abundance. The objective is to derive and implement the full chain
\[
P(k)\longrightarrow \sigma(M)\longrightarrow \nu(M,z)\longrightarrow f\longrightarrow \frac{\dd n}{\dd\ln M}
\]
without delegating the result to an external halo-mass-function package. External packages are reserved for independent validation. It records notation, units, derivations, numerical choices, and validation boundaries.
\end{abstract}

\section{Scope and convention freeze}

The later research asks whether dark-matter halo abundance retains information about an EDE episode \cite{PoulinEtAl2019,KamionkowskiRiess2023} beyond the instantaneous linear field and modern growth-history descriptions. That question is impossible to assess responsibly unless the mapping from linear fluctuations to conventional halo-abundance predictions is independently understood and numerically controlled.


Throughout the project, peak height means
\begin{equation}
\boxed{\nu(M,z)\equiv \frac{\delta_c(z)}{\sigma(M,z)}}.
\end{equation}
Some literature uses the symbol $\nu$ for $\delta_c^2/\sigma^2$. That alternative convention must never enter the source tree without an explicit conversion.

The h-scaled unit system used by the package is
\begin{align}
k &: h/\Mpc,\\
P(k) &: (\Mpc/h)^3,\\
R &: \Mpc/h,\\
M &: \Msun/h,\\
\rho &: (\Msun/h)/(\Mpc/h)^3,\\
\frac{\dd n}{\dd\ln M} &: (h/\Mpc)^3.
\end{align}

\section{Density contrast and the power spectrum}

Define the matter density contrast
\begin{equation}
\delta(\bm x,t)\equiv \frac{\rho(\bm x,t)-\bar\rho(t)}{\bar\rho(t)}.
\end{equation}
Using the Fourier convention
\begin{align}
\delta(\bm x) &= \int \frac{\dd^3k}{(2\pi)^3}\,\delta(\bm k)e^{i\bm k\cdot\bm x},\\
\delta(\bm k) &= \int \dd^3x\,\delta(\bm x)e^{-i\bm k\cdot\bm x},
\end{align}
statistical homogeneity and isotropy imply
\begin{equation}
\left\langle \delta(\bm k)\delta^*(\bm k')\right\rangle
=(2\pi)^3\delta_D^{(3)}(\bm k-\bm k')P(k).
\end{equation}
Because $\delta$ is dimensionless and $\delta_D^{(3)}(\bm k)$ has dimensions of inverse $k$-volume, $P(k)$ has dimensions of volume. The dimensionless power per logarithmic interval is
\begin{equation}
\boxed{\Delta^2(k)\equiv \frac{k^3P(k)}{2\pi^2}}.
\end{equation}

\section{Smoothing and the spherical top-hat window}

Halo abundance depends on fluctuations aggregated over a Lagrangian scale rather than the unsmoothed field at a point. Define
\begin{equation}
\delta_R(\bm x)=\int \dd^3x'\,W_R(|\bm x-\bm x'|)\delta(\bm x').
\end{equation}
For a unit-normalized real-space top hat,
\begin{equation}
W_R(r)=\begin{cases}
\dfrac{3}{4\pi R^3}, & r\le R,\\[4pt]
0, & r>R.
\end{cases}
\end{equation}
Its Fourier transform is
\begin{equation}
W(kR)=\frac{3}{R^3}\int_0^R r^2\frac{\sin(kr)}{kr}\dd r.
\end{equation}
Let $x=kR$ and $u=kr$. Then
\begin{align}
W(x)
&=\frac{3}{x^3}\int_0^x u\sin u\dd u\\
&=\frac{3}{x^3}\left[\sin x-x\cos x\right],
\end{align}
so
\begin{equation}
\boxed{W(x)=3\frac{\sin x-x\cos x}{x^3}}.
\end{equation}
Near $x=0$, direct evaluation is numerically ill-conditioned. Taylor expansion gives
\begin{equation}
\boxed{W(x)=1-\frac{x^2}{10}+\frac{x^4}{280}-\frac{x^6}{15120}+\mathcal O(x^8)}.
\end{equation}
The source implementation switches to this series locally to avoid catastrophic cancellation.

\section{Derivation of the smoothed variance}

The variance of the smoothed field is
\begin{equation}
\sigma^2(R)=\langle \delta_R^2(\bm x)\rangle.
\end{equation}
Since smoothing multiplies the Fourier amplitude by $W(kR)$,
\begin{equation}
\sigma^2(R)=\int\frac{\dd^3k}{(2\pi)^3}P(k)W^2(kR).
\end{equation}
Using isotropy, $\dd^3k=4\pi k^2\dd k$, hence
\begin{equation}
\boxed{\sigma^2(R)=\frac{1}{2\pi^2}\int_0^\infty k^2P(k)W^2(kR)\dd k}.
\end{equation}
Since $\dd k=k\dd\ln k$,
\begin{equation}
\boxed{\sigma^2(R)=\int_{-\infty}^{\infty}\Delta^2(k)W^2(kR)\dd\ln k}.
\end{equation}
The code integrates this logarithmic form. It is naturally matched to spectra spanning many decades in $k$ and makes the physical contribution per logarithmic interval explicit.

\section{Mass--radius mapping}

The Lagrangian mass inside a comoving top-hat radius $R$ is
\begin{equation}
\boxed{M=\frac{4\pi}{3}\rhobar R^3},
\end{equation}
with
\begin{equation}
\rhobar=\Omega_{m,0}\rho_{\mathrm{crit},0},
\qquad
\rho_{\mathrm{crit},0}=\frac{3H_0^2}{8\pi G}.
\end{equation}
In the h-scaled units used in the package,
\begin{equation}
\rho_{\mathrm{crit},0}=2.77536627\times10^{11}
\frac{\Msun/h}{(\Mpc/h)^3}.
\end{equation}
Thus
\begin{equation}
\boxed{R(M)=\left(\frac{3M}{4\pi\rhobar}\right)^{1/3}}.
\end{equation}
Combining this with the variance integral defines $\sigma(M)$.

\section{Linear growth}

For pressureless matter on linear sub-horizon scales,
\begin{equation}
\ddot D+2H\dot D-4\pi G\bar\rho_mD=0.
\end{equation}
Using $\dot D=aH D_a$ and differentiating once more,
\begin{equation}
\ddot D=a^2H^2D_{aa}+aH^2\left(1+\frac{aH_a}{H}\right)D_a.
\end{equation}
Substitution, division by $a^2H^2$, and $\bar\rho_m=3H_0^2\Omega_{m,0}a^{-3}/(8\pi G)$ give
\begin{equation}
\boxed{
D_{aa}+\left[\frac{3}{a}+\frac{\dd\ln H}{\dd a}\right]D_a
-\frac{3}{2}\frac{\Omega_{m,0}}{a^5E^2(a)}D=0,
}
\end{equation}
where $E(a)=H(a)/H_0$.

For a flat matter+Lambda reference background,
\begin{equation}
E^2(a)=\Omega_{m,0}a^{-3}+1-\Omega_{m,0}.
\end{equation}
The numerical solver starts sufficiently deep in matter domination with the growing-mode initial condition $D\propto a$, integrates with DOP853, and normalizes $D(1)=1$. The logarithmic growth rate is computed from the ODE state rather than by finite-differencing an interpolated solution:
\begin{equation}
\boxed{f(a)\equiv \frac{\dd\ln D}{\dd\ln a}=\frac{aD_a}{D}}.
\end{equation}
For Einstein--de Sitter, the exact growing mode is $D=a$ and $f=1$; these are hard validation tests.

\section{Peak height and collapse threshold}

The spherical-collapse reference used here is
\begin{equation}
\delta_c\simeq1.68647.
\end{equation}
Peak height is
\begin{equation}
\boxed{\nu(M,z)=\frac{\delta_c(z)}{\sigma(M,z)}}.
\end{equation}
Large $\nu$ corresponds to rare fluctuations. The later EDE project is particularly interested in high-redshift, high-$\nu$ halos because their abundance is exponentially sensitive to changes in the fluctuation statistics and collapse history.

\section{Press--Schechter derivation}

Press and Schechter introduced the foundational analytic description of halo abundance \cite{PressSchechter1974}. For a Gaussian smoothed overdensity,
\begin{equation}
p(\delta;M)=\frac{1}{\sqrt{2\pi}\sigma(M)}
\exp\left[-\frac{\delta^2}{2\sigma^2(M)}\right].
\end{equation}
The probability above a constant collapse barrier is
\begin{equation}
P(\delta>\delta_c)=\frac{1}{2}\operatorname{erfc}\left(\frac{\delta_c}{\sqrt{2}\sigma}\right).
\end{equation}
The original Press--Schechter mass accounting introduces a factor of two, yielding
\begin{equation}
F(>M)=\operatorname{erfc}\left(\frac{\nu}{\sqrt2}\right).
\end{equation}
Since
\begin{equation}
\frac{\dd}{\dd x}\operatorname{erfc}(x)=-\frac{2}{\sqrt\pi}e^{-x^2},
\end{equation}
we obtain
\begin{equation}
\frac{\dd F}{\dd\nu}=-\sqrt{\frac{2}{\pi}}e^{-\nu^2/2}.
\end{equation}
The number density follows from converting mass fraction to objects per mass interval:
\begin{equation}
\frac{\dd n}{\dd M}=-\frac{\rhobar}{M}\frac{\dd F}{\dd M}.
\end{equation}
With $\nu=\delta_c/\sigma$,
\begin{equation}
\frac{\dd\nu}{\dd M}=\nu\frac{\dd\ln\sigma^{-1}}{\dd M},
\end{equation}
therefore
\begin{equation}
\boxed{
\frac{\dd n}{\dd M}=\frac{\rhobar}{M}
\sqrt{\frac{2}{\pi}}\nu e^{-\nu^2/2}
\frac{\dd\ln\sigma^{-1}}{\dd M}.
}
\end{equation}
Define the project-convention multiplicity
\begin{equation}
\boxed{f_{\rm PS}(\nu)=\sqrt{\frac{2}{\pi}}\nu e^{-\nu^2/2}},
\end{equation}
so
\begin{equation}
\boxed{
\frac{\dd n}{\dd\ln M}=\frac{\rhobar}{M}f(\nu)
\frac{\dd\ln\sigma^{-1}}{\dd\ln M}.
}
\end{equation}
This generic transformation is implemented exactly once in the source tree.

\subsection{Mass-fraction normalization}

Under the adopted convention, $f(\nu)\dd\ln\nu$ is a differential mass fraction. For Press--Schechter,
\begin{align}
\int_0^\infty f_{\rm PS}(\nu)\dd\ln\nu
&=\sqrt{\frac{2}{\pi}}\int_0^\infty e^{-\nu^2/2}\dd\nu\\
&=1.
\end{align}
This identity is tested numerically to high accuracy.

\section{Excursion sets and the cloud-in-cloud problem}

The factor of two in the original argument signals a deeper counting problem: a point that appears uncollapsed on a small smoothing scale can lie inside a larger collapsed region. Bond et~al.\ \cite{BondEtAl1991} formulate hierarchical collapse as a first-crossing problem.

Let
\begin{equation}
S\equiv\sigma^2.
\end{equation}
For a sharp-$k$ filter, successive increments form a Markov random walk and the density of trajectories obeys
\begin{equation}
\frac{\partial\Pi}{\partial S}=\frac{1}{2}\frac{\partial^2\Pi}{\partial\delta^2}
\end{equation}
with absorbing barrier $\Pi(\delta_c,S)=0$. The image solution below the barrier is
\begin{equation}
\Pi(\delta,S)=\frac{1}{\sqrt{2\pi S}}
\left[
 e^{-\delta^2/(2S)}-e^{-(2\delta_c-\delta)^2/(2S)}
\right].
\end{equation}
The first-crossing distribution is
\begin{equation}
\boxed{
\mathcal F(S)=\frac{\delta_c}{\sqrt{2\pi}}S^{-3/2}
\exp\left(-\frac{\delta_c^2}{2S}\right).
}
\end{equation}
Using $S=\sigma^2$ and
\begin{equation}
\frac{\dd n}{\dd M}=\frac{\rhobar}{M}\mathcal F(S)\left|\frac{\dd S}{\dd M}\right|,
\end{equation}
recovers the Press--Schechter mass function without inserting an ad hoc mass-doubling factor. The package does not yet implement a Monte Carlo excursion-set solver, but this derivation explains its normalization tests.

\section{Sheth--Tormen and ellipsoidal collapse}

Simulation-calibrated departures from spherical Press--Schechter behavior motivated the Sheth--Tormen form \cite{ShethTormen1999}, with ellipsoidal/moving-barrier interpretation developed further by Sheth, Mo and Tormen \cite{ShethMoTormen2001}. In the project's $\nu=\delta_c/\sigma$ convention,
\begin{equation}
\boxed{
f_{\rm ST}(\nu)=A\sqrt{\frac{2a}{\pi}}
\left[1+(a\nu^2)^{-p}\right]\nu
\exp\left(-\frac{a\nu^2}{2}\right).
}
\end{equation}
The common parameter choice is $a=0.707$ and $p=0.3$. Rather than treating $A$ as an unexplained constant, impose
\begin{equation}
\int_0^\infty f_{\rm ST}(\nu)\dd\ln\nu=1.
\end{equation}
The first term contributes $A$ and the correction term contributes
\begin{equation}
A\frac{2^{-p}}{\sqrt\pi}\Gamma\left(\frac12-p\right),
\end{equation}
therefore
\begin{equation}
\boxed{
A^{-1}=1+\frac{2^{-p}}{\sqrt\pi}\Gamma\left(\frac12-p\right).
}
\end{equation}
For $p=0.3$, $A\simeq0.32218$. A particularly strong regression test is the analytic limit $a=1$, $p=0$. In that case $A=1/2$ and the Sheth--Tormen expression reduces exactly to Press--Schechter.

\section{Tinker et~al.\ and the limits of universality}

Tinker et~al.\ \cite{TinkerEtAl2008} calibrate an empirical spherical-overdensity HMF directly from simulations. Their fitting form is
\begin{equation}
\boxed{
f(\sigma)=A\left[\left(\frac{\sigma}{b}\right)^{-a}+1\right]
\exp\left(-\frac{c}{\sigma^2}\right).
}
\end{equation}
For the $\Delta=200$ relative-to-mean, $z=0$ reference calibration, the implementation uses
\begin{equation}
A=0.186,\qquad a=1.47,\qquad b=2.57,\qquad c=1.19.
\end{equation}
The restricted class name in the code intentionally records the calibration domain. General redshift and overdensity interpolation is postponed until the halo-definition machinery exists. Tinker et~al.\ also show that the HMF is not exactly universal at precision level; this is an important precursor to the later EDE question.

\section{Why Evolution Mapping III changes the later research question}

Fiorilli et~al.\ \cite{FiorilliEtAl2026} model HMF non-universality using recent structure-formation history together with local linear-power-spectrum shape, reaching percent-level accuracy over their tested domain. Consequently, the later EDE project cannot claim novelty merely by showing that an older universal fit has residuals. It must test whether EDE leaves information beyond modern growth-history and spectrum descriptors.

\section{Numerical design and code architecture}

\subsection{No silent power-spectrum extrapolation}
The tabulated power-spectrum object validates positivity and monotonic $k$ ordering and rejects out-of-domain evaluation by default. Variance integration uses only supplied support. Each $\sigma(R)$ result carries low- and high-$k$ edge ratios: the integrand value at each boundary divided by its maximum. These diagnostics identify mass scales for which the integral may be artificially truncated.

\subsection{Shape-preserving and analytical HMF Jacobians}
The HMF requires
\begin{equation}
J(M)=\frac{\dd\ln\sigma^{-1}}{\dd\ln M}.
\end{equation}
The source computes $J$ either by differentiating a shape-preserving PCHIP interpolant of $\ln\sigma$ versus $\ln M$ \cite{FritschButland1984}, or by evaluating the exact analytical logarithmic derivative:
\begin{equation}
J(M) = -\frac{R}{6\sigma^2}\frac{\dd\sigma^2}{\dd R} = -\frac{1}{3\sigma^2}\int\dd\ln k\,\Delta^2(k)W(kR)\left[kR\,W'(kR)\right]
\end{equation}
directly from the analytical top-hat window derivative $W'(x) = \frac{3}{x}(\frac{\sin x}{x} - W(x))$. The exact same Jacobian is shared by every multiplicity model, preventing derivative implementation differences from masquerading as model differences.

\subsection{Growth derivative from the ODE state}
The numerical ODE state contains both $D$ and $D_a$. The code computes $f=aD_a/D$ directly. Differentiating a second interpolant of $D(a)$ would add an unnecessary numerical operation and another source of error.

\subsection{External implementations are validators}
CLASS \cite{Blas2011}, Colossus \cite{Diemer2018}, and the Python \texttt{hmf} package \cite{Murray2013} are deliberately not imported by the scientific core. They are used to compare $P(k)$, $\sigma_8$, $\sigma(M)$ and HMF outputs under matched conventions. Independent agreement is evidence; self-agreement through a wrapped external implementation is not.


\section{foundational validation program}

The hard tests are:
\begin{enumerate}
\item $M\rightarrow R\rightarrow M$ round trip to floating-point precision.
\item Einstein--de Sitter exact growing mode $D(a)=a$.
\item Einstein--de Sitter growth rate $f=1$.
\item Power scaling $P\rightarrow4P$ implies $\sigma\rightarrow2\sigma$.
\item Press--Schechter $\int f\dd\ln\nu=1$.
\item Sheth--Tormen $\int f\dd\ln\nu=1$.
\item Sheth--Tormen $a=1,p=0$ reduces exactly to Press--Schechter.
\item A synthetic power law $\sigma\propto M^{-\alpha}$ yields $\dd\ln\sigma^{-1}/\dd\ln M=\alpha$.
\item Every HMF output is finite and non-negative.
\item Variance conclusions are rejected when the finite-$k$ edge diagnostic is not acceptably small.
\end{enumerate}

\section{Questions this foundation must answer}

Without code or notes, the researcher must be able to:
\begin{enumerate}
\item derive the top-hat Fourier transform;
\item derive the variance integral from the power-spectrum two-point function;
\item explain every unit in the chain $P(k)\rightarrow\dd n/\dd\ln M$;
\item derive Press--Schechter and explain the original factor-of-two problem;
\item derive the Bond et~al.\ first-crossing distribution at the level shown here;
\item explain why Sheth--Tormen changes both low- and high-$\nu$ behavior;
\item explain why a Tinker fit is inseparable from its mass definition and calibration domain;
\item explain why the 2026 growth-history HMF work raises the novelty bar for EDE;
\item identify at least three numerical ways a correct-looking HMF curve can nevertheless be wrong.
\end{enumerate}

The final conceptual standard is simple: if the independent $P(k)\rightarrow\sigma(M)\rightarrow\mathrm{HMF}$ layer is wrong, every later EDE comparison and scientific claim can be consistently wrong. This foundation therefore prioritizes auditable correctness before novelty.

\bibliographystyle{unsrt}
\bibliography{references}

\end{document}
